\documentclass[11pt]{article}
\usepackage[margin=1in]{geometry}
\usepackage{amsmath,amssymb,amsthm,mathtools}
\usepackage{booktabs}
\usepackage{array}
\usepackage{enumitem}
\usepackage{hyperref}
\usepackage{microtype}
\usepackage[T1]{fontenc}

\hypersetup{colorlinks=true,linkcolor=blue,citecolor=blue,urlcolor=blue}

\newtheorem{theorem}{Theorem}[section]
\newtheorem{proposition}[theorem]{Proposition}
\newtheorem{lemma}[theorem]{Lemma}
\newtheorem{corollary}[theorem]{Corollary}
\newtheorem{conjecture}[theorem]{Conjecture}
\theoremstyle{remark}
\newtheorem{remark}[theorem]{Remark}

\newcommand{\vTwo}{\nu_2}
\newcommand{\Nodd}{\mathbb N_{\mathrm{odd}}}
\newcommand{\eps}{\varepsilon}

\title{\textbf{A $5^r$-Scaled Family of the Reduced Ternary $(5k\pm1)/4$ Collatz-Type Map}\\
\large Exact Algebraic Conjugacy, Cycle Preservation, Orbit-Statistic Invariance, and Transfer of the $10^{10}$ Finite Verification}
\author{Farhad Banazadeh\\
\texttt{fn.bana@hotmail.com}\\
ORCID: 0009-0004-7023-0298}
\date{11 September 2026}

\begin{document}
\maketitle

\begin{abstract}
We develop the full $5^r$-scaled family associated with the residue-dependent reduced ternary $(5k\pm1)/4$ Collatz-type map studied in the author's preceding work.  The base state space is the positive odd integers and the base affine numerator is
\[
M_0(y)=
\begin{cases}
5y-1,&y\equiv1\pmod4,\\
5y+1,&y\equiv3\pmod4,
\end{cases}
\qquad
T_0(y)=\frac{M_0(y)}{2^{\vTwo(M_0(y))}}.
\]
For every integer $r\ge0$ we define the scaled domain $S_r=5^r\Nodd$ and
\[
M_r(x)=
\begin{cases}
5x-5^r,&x\equiv1\pmod4,\\
5x+5^r,&x\equiv3\pmod4,
\end{cases}
\qquad
T_r(x)=\frac{M_r(x)}{2^{\vTwo(M_r(x))}}.
\]
Thus the unreduced branches have the requested form $(5k\pm5^r)/4$ after the forced factor $4$ is removed, followed, as in the base system, by complete removal of any additional powers of $2$.

The central result is an exact algebraic conjugacy.  Since $5^r\equiv1\pmod4$ and $\gcd(5^r,2)=1$, multiplication by $5^r$ preserves both the branch residue and the complete two-adic valuation.  Consequently
\[
T_r(5^r y)=5^rT_0(y),\qquad
T_r^j(5^r y)=5^rT_0^j(y)\quad(j\ge0).
\]
Hence every orbit relation, branch sequence, valuation sequence, periodic orbit, minimal cycle length, capture time, first-descent stopping time, basin membership, and normalized orbit geometry is transported bijectively between the base system and its scaled copy.  Absolute reduced and raw affine peaks are multiplied by $5^r$, while peak/start ratios are unchanged.  In particular, the base fixed point $1$ becomes $5^r$, and the base three-cycle $7\to9\to11\to7$ becomes
\[
7\cdot5^r\to9\cdot5^r\to11\cdot5^r\to7\cdot5^r.
\]
No additional periodic orbit can occur inside $S_r$ unless a corresponding base periodic orbit exists.

The preceding base paper exhaustively verified all $5,000,000,000$ positive odd starts below $10^{10}$.  We do not relabel that computation as a new scan.  Instead, the conjugacy theorem transfers each verified base trajectory exactly to the corresponding scaled set
\[
\{5^r y:1\le y<10^{10},\ y\ {\rm odd}\}\subset S_r,
\]
which reaches absolute scale below $5^r10^{10}$.  The global two-attractor convergence statement remains open: the present theorem proves exact equivalence of the scaled and base dynamical problems, not global convergence of the unresolved base problem.
\end{abstract}

\noindent\textbf{Keywords:} Collatz-type map; $5x\pm1$; $(5k\pm5^r)/4$; $5^r$ scaling; algebraic conjugacy; two-adic valuation; periodic orbits; cycle preservation; computational number theory.

\section{Introduction}

The reduced $5x\pm1$ map considered in the author's base paper is a residue-dependent odd-to-odd dynamical system.  For an odd state $y$, the sign is chosen from $y\bmod4$ so that $5y\pm1$ is divisible by at least $4$, after which all powers of $2$ are removed.  The base study developed closure, an exact periodic-orbit identity, a two-adic valuation law and drift heuristic, inverse branches and basin structure, and an exhaustive finite verification through $10^{10}$.

The present paper asks a structural rather than a new billion-scale computational question.  Replace the constants $\pm1$ by $\pm5^r$ and simultaneously restrict the state space to the $5^r$-scaled copy of the base odd domain.  Does this create genuinely new dynamics?  We prove that it does not.  For every integer $r\ge0$, the resulting system is exactly conjugate to the base map through the linear scaling
\[
\Phi_r(y)=5^r y.
\]

This is stronger than numerical similarity for selected $r$.  Exact conjugacy proves the correspondence for every $r\ge0$ at once and identifies precisely what is preserved: full trajectories, branch choices, complete two-adic valuations, cycle lengths, transient/capture times, first-descent stopping times, basin membership, and all normalized orbit ratios.  Only dimensional quantities such as state values and peaks are multiplied by $5^r$.

The finite computation in the base article remains base-system evidence.  Its numerical counts are inherited here through a theorem and are not presented as fresh scans of infinitely many scaled systems.  This distinction is essential for reproducibility and for separating exact algebraic consequences from finite computational evidence.

\section{The base reduced ternary map}

Let
\[
S=\Nodd=\{1,3,5,\ldots\}.
\]
Define the sign selector
\[
\eps(y)=
\begin{cases}
-1,&y\equiv1\pmod4,\\
+1,&y\equiv3\pmod4.
\end{cases}
\]
The base affine numerator and reduced map are
\begin{equation}
M_0(y)=5y+\eps(y),
\qquad
T_0(y)=\frac{M_0(y)}{2^{\vTwo(M_0(y))}}.
\label{eq:base}
\end{equation}
Because $5\equiv1\pmod4$, one has $4\mid M_0(y)$ for every odd $y$, so $\vTwo(M_0(y))\ge2$ and $T_0(y)$ is again odd.

Two exact periodic attractors displayed in the base work are
\[
1\to1
\]
and
\[
7\to9\to11\to7,
\]
since $5\cdot7+1=36=4\cdot9$, $5\cdot9-1=44=4\cdot11$, and $5\cdot11+1=56=8\cdot7$.

\section{The $5^r$-scaled family}

Fix an integer $r\ge0$.

\subsection{Scaled domain}

Define
\begin{equation}
S_r=5^rS=\{5^r y:y\in S\}.
\label{eq:domain}
\end{equation}
Every $x\in S_r$ has a unique representation $x=5^r y$ with $y$ odd.  Since
\[
5^r\equiv1\pmod4,
\]
we have
\begin{equation}
5^r y\equiv y\pmod4.
\label{eq:residuepreserve}
\end{equation}
Thus the scaling preserves the branch selector exactly:
\[
\eps(5^r y)=\eps(y).
\]

\subsection{Scaled affine and reduced maps}

For $x\in S_r$, define
\begin{equation}
M_r(x)=5x+5^r\eps(x)
=
\begin{cases}
5x-5^r,&x\equiv1\pmod4,\\
5x+5^r,&x\equiv3\pmod4.
\end{cases}
\label{eq:scaledaffine}
\end{equation}
The corresponding completely reduced map is
\begin{equation}
T_r(x)=\frac{M_r(x)}{2^{\vTwo(M_r(x))}}.
\label{eq:scaledmap}
\end{equation}

Because $5^r\equiv1\pmod4$, the branch rule gives $M_r(x)\equiv0\pmod4$ for every $x\in S_r$.  Therefore the commonly written first compressed form is
\[
\frac{5x\pm5^r}{4},
\]
but the dynamical map \eqref{eq:scaledmap} continues to remove every additional factor of $2$, exactly as in the base paper.

\begin{remark}[Why the scaled domain matters]
The exact theorem concerns $S_r=5^rS$.  For a general odd integer $x$ not divisible by $5^r$, the factorization used below is unavailable.  Therefore ``the same dynamics'' is asserted on the scaled copy $S_r$, not on all odd integers under the formula $5x\pm5^r$.
\end{remark}

\section{The key two-adic scaling identities}

\begin{lemma}[Affine scaling]
For every $r\ge0$ and $y\in S$,
\begin{equation}
M_r(5^r y)=5^rM_0(y).
\label{eq:affinescale}
\end{equation}
\end{lemma}

\begin{proof}
By \eqref{eq:residuepreserve}, $\eps(5^r y)=\eps(y)$.  Hence
\[
M_r(5^r y)
=5(5^r y)+5^r\eps(y)
=5^r(5y+\eps(y))
=5^rM_0(y).
\]
\end{proof}

\begin{lemma}[Valuation invariance]
For every nonzero integer $n$ and every $r\ge0$,
\begin{equation}
\vTwo(5^r n)=\vTwo(n).
\label{eq:valscale}
\end{equation}
\end{lemma}

\begin{proof}
The factor $5^r$ is odd and therefore contributes no factor of $2$.
\end{proof}

Combining the lemmas,
\begin{equation}
\vTwo(M_r(5^r y))=\vTwo(M_0(y)).
\label{eq:matchedvaluation}
\end{equation}
Thus corresponding trajectories do not merely scale as sets of states; they remove exactly the same power of $2$ at every matched step.

\section{Exact algebraic conjugacy}

Define the scaling bijection
\[
\Phi_r:S\to S_r,\qquad \Phi_r(y)=5^r y,
\]
with inverse $\Phi_r^{-1}(x)=x/5^r$.

\begin{theorem}[One-step conjugacy]
For every $r\ge0$ and $y\in S$,
\begin{equation}
T_r(5^r y)=5^rT_0(y).
\label{eq:onestep}
\end{equation}
Equivalently,
\[
T_r\circ\Phi_r=\Phi_r\circ T_0,
\qquad
T_r=\Phi_r\circ T_0\circ\Phi_r^{-1}.
\]
\end{theorem}

\begin{proof}
Using \eqref{eq:affinescale} and \eqref{eq:matchedvaluation},
\[
T_r(5^r y)
=\frac{5^rM_0(y)}{2^{\vTwo(M_0(y))}}
=5^rT_0(y).
\]
\end{proof}

\begin{theorem}[All iterates]
For every $j\ge0$, $r\ge0$, and $y\in S$,
\begin{equation}
T_r^j(5^r y)=5^rT_0^j(y).
\label{eq:alliterates}
\end{equation}
\end{theorem}

\begin{proof}
The case $j=0$ is immediate.  If the identity holds for $j$, then
\[
T_r^{j+1}(5^r y)
=T_r(5^rT_0^j(y))
=5^rT_0(T_0^j(y))
=5^rT_0^{j+1}(y).
\]
\end{proof}

Equation \eqref{eq:alliterates} is the organizing identity of the paper.  Every statement below follows either immediately from it or from the stepwise valuation identity \eqref{eq:matchedvaluation}.

\section{Exact preservation of periodic orbits}

\begin{theorem}[Cycle bijection]
A tuple $C=(c_0,\ldots,c_{L-1})$ is a cycle of $T_0$ if and only if
\[
5^rC=(5^rc_0,\ldots,5^rc_{L-1})
\]
is a cycle of $T_r$.  Minimal cycle length is preserved.
\end{theorem}

\begin{proof}
If $T_0^L(c_0)=c_0$, then
\[
T_r^L(5^rc_0)=5^rT_0^L(c_0)=5^rc_0.
\]
Conversely, if $x=5^ry\in S_r$ and $T_r^L(x)=x$, division of
\[
5^rT_0^L(y)=5^ry
\]
by $5^r$ gives $T_0^L(y)=y$.  Bijectivity of $\Phi_r$ preserves minimal period.
\end{proof}

\begin{corollary}
No periodic orbit exists in $S_r$ that is not the $5^r$-multiple of a base periodic orbit.  Scaling introduces no new periodic structure inside the scaled domain.
\end{corollary}

The displayed base attractors become
\begin{equation}
5^r\to5^r
\label{eq:fixedscaled}
\end{equation}
and
\begin{equation}
7\cdot5^r\to9\cdot5^r\to11\cdot5^r\to7\cdot5^r.
\label{eq:cyclescaled}
\end{equation}
For example,
\[
r=1:\quad35\to45\to55\to35,
\]
\[
r=2:\quad175\to225\to275\to175,
\]
\[
r=3:\quad875\to1125\to1375\to875.
\]
These are theorem-level consequences, not independent numerical discoveries.

\section{The periodic-orbit equation under scaling}

For a base orbit write
\[
2^{a_i}y_{i+1}=5y_i+\eps_i,
\qquad
a_i=\vTwo(5y_i+\eps_i)\ge2.
\]
For the scaled orbit $x_i=5^ry_i$,
\begin{equation}
2^{a_i}x_{i+1}=5x_i+5^r\eps_i.
\label{eq:scaledstepcycle}
\end{equation}
Dividing by $5^r$ recovers the base relation exactly.

Let
\[
A=\sum_{i=0}^{L-1}a_i,\qquad
S_j=\sum_{i=0}^{j-1}a_i,\quad S_0=0.
\]
Iteration of \eqref{eq:scaledstepcycle} around a length-$L$ cycle gives
\begin{equation}
(2^A-5^L)x_0
=
5^r\sum_{j=0}^{L-1}5^{L-1-j}\eps_j2^{S_j}.
\label{eq:scaledcycleidentity}
\end{equation}
Hence
\begin{equation}
x_0
=
5^r
\frac{\displaystyle\sum_{j=0}^{L-1}5^{L-1-j}\eps_j2^{S_j}}
{\displaystyle 2^A-5^L}
=5^ry_0.
\label{eq:scaledcycleclosed}
\end{equation}
Thus the scaled Diophantine cycle equation contains no new integrality condition beyond the base equation.  This supplies a second direct proof of cycle preservation.

\section{Orbit statistics preserved by conjugacy}

Let
\[
y_j=T_0^j(y),\qquad x_j=T_r^j(5^ry).
\]
Then $x_j=5^ry_j$ for all $j$.

\subsection{Capture/transient time}

Let $\tau_0(y)$ denote the first iteration at which the base orbit enters one of the displayed attractors $\{1\}$ or $\{7,9,11\}$.  Then
\begin{equation}
\tau_r(5^ry)=\tau_0(y).
\label{eq:capturetime}
\end{equation}
The complete valuation sequence $(a_0,a_1,\ldots)$ is identical, so the accumulated number of individual divisions by $2$ over any matched finite segment is also identical.

\subsection{First-descent stopping time}

Independently of capture time, define the first-descent stopping time
\[
\sigma_0(y)=\min\{j\ge1:T_0^j(y)<y\},
\]
whenever such a $j$ exists.  Since multiplication by $5^r$ preserves strict inequalities,
\[
T_r^j(5^ry)<5^ry
\iff
T_0^j(y)<y.
\]
Therefore
\begin{equation}
\sigma_r(5^ry)=\sigma_0(y).
\label{eq:stopping}
\end{equation}
If
\[
L_0(y)=T_0^{\sigma_0(y)}(y)
\]
is the first lower value, then
\begin{equation}
L_r(5^ry)=5^rL_0(y).
\label{eq:firstlower}
\end{equation}

\subsection{Peaks and normalized ratios}

For any matched finite trajectory segment,
\begin{equation}
\max_j x_j=5^r\max_j y_j.
\label{eq:peak}
\end{equation}
The raw affine values scale in the same way because $M_r(5^ry)=5^rM_0(y)$.  Thus both reduced and raw affine peaks are multiplied by $5^r$.  On the other hand,
\begin{equation}
\frac{x_j}{x_0}=\frac{y_j}{y_0},
\label{eq:ratio}
\end{equation}
so every normalized peak/start ratio and every normalized trajectory profile is invariant.  The iteration index at which a matched peak occurs is also unchanged.

\subsection{Basin membership}

For a base cycle $C$, let $B_0(C)$ be its basin and let $B_r(5^rC)$ be the basin of the scaled cycle within $S_r$.  Then
\begin{equation}
5^ry\in B_r(5^rC)
\iff
y\in B_0(C).
\label{eq:basin}
\end{equation}
Consequently basin counts and basin proportions are exactly preserved on corresponding finite sets.  They should not be compared at the same absolute cutoff without accounting for the different scaled domains.

\section{Two-adic valuation law and drift under scaling}

The base map forces $a=\vTwo(M_0(y))\ge2$.  In the natural two-adic residue model,
\[
\Pr(a=k)=2^{-(k-1)},\qquad k\ge2,
\]
and
\[
\mathbb E[a]=3.
\]
For large $y$, the base logarithmic increment is approximated by
\[
\log\frac{T_0(y)}{y}\approx\log5-a\log2,
\]
giving
\begin{equation}
\mathbb E[\Delta\log]\approx\log5-3\log2=\log(5/8)\approx-0.4700036292.
\label{eq:drift}
\end{equation}
The geometric typical multiplier is therefore $5/8$.  The ordinary expected approximate multiplier is instead
\[
\mathbb E\!\left[\frac5{2^a}\right]=\frac56.
\]

For matched states the normalized one-step ratio is exact:
\begin{equation}
\frac{T_r(5^ry)}{5^ry}=\frac{T_0(y)}{y}.
\label{eq:normalizedstep}
\end{equation}
Hence the scaled family inherits the same valuation distribution, the same normalized multiplicative process, the same critical average valuation $\log_2 5$, and the same drift heuristic.  As in the base paper, negative expected logarithmic drift is a probabilistic interpretation and not a proof that every orbit converges.

\section{Exact inverse branches and scaled inverse forests}

Suppose $T_0(y)=z$ and the removed power is $2^a$, $a\ge2$.  The base inverse candidates are
\[
y=\frac{2^az+1}{5}
\quad\text{or}\quad
y=\frac{2^az-1}{5},
\]
subject to integrality and the corresponding branch condition.

For the scaled system, if $T_r(x)=w$ with $x=5^ry$ and $w=5^rz$, then
\[
x=5^r\frac{2^az+1}{5}
=\frac{2^aw+5^r}{5},
\]
or
\[
x=5^r\frac{2^az-1}{5}
=\frac{2^aw-5^r}{5}.
\]
Thus every inverse edge in the base forest is multiplied by $5^r$, and conversely division by $5^r$ sends every inverse edge in $S_r$ back to a base edge.  The entire rooted inverse forest of the fixed point and the entire rooted inverse forest of the three-cycle are therefore scaled copies, not merely forward-orbit copies.

Because the inverse-tree geometry is transported bijectively, any basin-density statement formulated intrinsically on corresponding scaled sets is equivalent to its base statement.  The observed finite $72/28$-type split is inherited on corresponding finite sets; no new limiting-density theorem follows from scaling alone.

\section{Finite verification inherited from the base paper}

The base computation exhaustively tested all positive odd starting values below $10^{10}$:
\[
5,000,000,000\ \text{starts}.
\]
Every tested orbit reached one of the two displayed attractors within the stated safety cap.  The inherited base summary is:
\begin{center}
\begin{tabular}{@{}lr@{}}
\toprule
Quantity & Base value through $10^{10}$\\
\midrule
Odd starts & $5,000,000,000$\\
Basin of $1$ & $3,599,424,753$ ($71.988495060\%$)\\
Basin of $7\to9\to11\to7$ & $1,400,575,247$ ($28.011504940\%$)\\
Unresolved & $0$\\
Mean capture time & $43.775639881800$\\
Maximum capture time & $174$ at $8,793,787,045$\\
Maximum reduced odd peak & $11,134,926,214,451$\\
Peak-holder start & $9,250,534,785$\\
Maximum raw affine peak & $55,674,631,072,256$\\
Maximum observed $\vTwo$ & $34$\\
Maximum reduced-peak/start ratio & $3673.419845666274$\\
\bottomrule
\end{tabular}
\end{center}

These are not new measurements of the $5^r$ family.  The exact theorem implies that for every fixed $r$, the corresponding set
\begin{equation}
F_r=\{5^ry:1\le y<10^{10},\ y\ {\rm odd}\}
\label{eq:finitefamily}
\end{equation}
contains exactly $5,000,000,000$ scaled starts and inherits the same basin counts, proportions, capture-time distribution, maximum capture time, and valuation sequences.

Absolute values scale.  In particular, on $F_r$:
\[
\text{maximum reduced peak}
=
5^r\cdot11,134,926,214,451,
\]
attained by
\[
5^r\cdot9,250,534,785,
\]
and the maximum raw affine peak is
\[
5^r\cdot55,674,631,072,256.
\]
The maximum capture time remains $174$, attained at the scaled start
\[
5^r\cdot8,793,787,045.
\]
The maximum observed valuation remains $34$, and the maximum normalized reduced-peak/start ratio remains exactly the same numerical ratio, $3673.419845666274$.

The corresponding finite set lies below absolute scale $5^r10^{10}$, but it is a structured subset $S_r$, not the set of all odd integers below that absolute bound.

\section{Finite transfer theorem and global equivalence}

\begin{theorem}[Finite transfer theorem]
For every $r\ge0$, every scaled start
\[
x=5^ry,\qquad 1\le y<10^{10},\quad y\ {\rm odd},
\]
enters either the fixed point $5^r$ or the scaled three-cycle
\[
7\cdot5^r\to9\cdot5^r\to11\cdot5^r\to7\cdot5^r
\]
under $T_r$, subject exactly to the same finite computational conditions used in the base verification.
\end{theorem}

\begin{proof}
The base computation establishes the finite statement for each corresponding $y$.  The all-iterates identity \eqref{eq:alliterates} transports the entire verified trajectory by multiplication by $5^r$.
\end{proof}

\begin{conjecture}[Scaled two-attractor convergence]
For every $r\ge0$ and every $x\in S_r$, the orbit of $x$ under $T_r$ eventually enters either $5^r$ or the scaled three-cycle \eqref{eq:cyclescaled}.
\end{conjecture}

\begin{theorem}[Equivalence of global conjectures]
The scaled two-attractor conjecture holds for any fixed $r\ge0$ if and only if the base two-attractor conjecture holds on $S$.  Consequently it holds for all $r\ge0$ if and only if it holds for $r=0$.
\end{theorem}

\begin{proof}
The conjugacy $\Phi_r$ transports every orbit in both directions.
\end{proof}

Thus the parameter $r$ creates no independent global convergence problem.  It also does not solve the base problem.  A divergent base orbit, if one exists, scales to a divergent orbit in every $S_r$; an undiscovered base cycle scales to a cycle of the same length in every $S_r$.

\section{What is proved and what is not}

The exact algebraic results prove, for every $r\ge0$:
\begin{enumerate}[label=(\roman*)]
\item $S_r=5^rS$ is in bijection with the base domain;
\item multiplication by $5^r$ preserves the relevant residue modulo $4$;
\item $M_r(5^ry)=5^rM_0(y)$;
\item the complete two-adic valuation is unchanged at corresponding steps;
\item $T_r$ and $T_0$ are exactly conjugate;
\item all iterates correspond exactly;
\item branch sequences and valuation sequences correspond exactly;
\item periodic orbits and minimal cycle lengths correspond bijectively;
\item capture/transient times are preserved;
\item first-descent stopping times are preserved and first-lower values scale by $5^r$;
\item absolute reduced and raw peaks scale by $5^r$;
\item peak locations in iteration time and all normalized peak/start ratios are preserved;
\item basin membership and basin counts on corresponding finite sets transfer exactly;
\item inverse branches and the complete inverse forests scale exactly;
\item the base finite verification through $10^{10}$ transfers to the corresponding scaled set through absolute scale $5^r10^{10}$.
\end{enumerate}

What is \emph{not} proved is global convergence of every positive base orbit, nonexistence of additional base cycles, or existence and exact values of limiting basin densities.  A finite scan through $10^{10}$ cannot exclude phenomena beyond the tested range.  Since the scaled systems are conjugate to the base system, they inherit exactly this unresolved global status.

\section{Reproducibility and computational verification of the identity}

The supplementary archive accompanying this paper contains the \LaTeX{} source, compiled PDF, a small Python verifier, a machine-readable JSON statement of the scaling rules and inherited base statistics, a README, the inherited base summary, and SHA-256 checksums.

The verifier checks, for user-selected finite ranges of $r$ and odd $y$, the identities
\[
M_r(5^ry)=5^rM_0(y),
\]
\[
\vTwo(M_r(5^ry))=\vTwo(M_0(y)),
\]
\[
T_r(5^ry)=5^rT_0(y),
\]
and compares several iterates.  This verifier is a consistency check, not the proof.  The proof is the algebraic factorization together with the facts $5^r\equiv1\pmod4$ and $\gcd(5^r,2)=1$.

The $10^{10}$ numerical evidence belongs to the base record.  The present paper deliberately avoids presenting inherited trajectories as newly recomputed data.

\section{Conclusion}

Replacing the base constants $\pm1$ by $\pm5^r$ while scaling the odd state space by $5^r$ produces an exact scaled copy of the reduced ternary $(5k\pm1)/4$ dynamics.  The generalized affine branches are
\[
5x-5^r,\qquad 5x+5^r,
\]
and their forced first compression is of the form $(5x\pm5^r)/4$.  After complete two-adic reduction,
\[
T_r(5^ry)=5^rT_0(y)
\]
for every $r\ge0$ and every positive odd $y$.

The reason is exact and elementary: powers of $5$ are simultaneously congruent to $1$ modulo $4$ and coprime to $2$.  The first fact preserves branch selection; the second preserves the complete two-adic valuation.  Consequently every scaled trajectory is precisely the $5^r$-multiple of one base trajectory.

Every scaled periodic orbit is the $5^r$-multiple of a base periodic orbit with the same minimal period.  The fixed point and the three-cycle therefore become $5^r$ and
\[
7\cdot5^r\to9\cdot5^r\to11\cdot5^r\to7\cdot5^r.
\]
Capture times and first-descent stopping times are unchanged; first-lower values and absolute peaks scale by $5^r$; normalized peak ratios are unchanged; and the inverse forests and basin memberships are transported bijectively.

The exhaustive base verification of all five billion odd starts below $10^{10}$ therefore transfers, without a new billion-scale scan, to every corresponding $5^r$-scaled finite domain.  At the same time, the global convergence question is unchanged.  Proving the two-attractor conjecture for the base system would prove it simultaneously for the whole scaled family, while any base counterexample would propagate throughout the family.  The parameter $r$ is therefore an exact scaling symmetry, not a source of independent dynamics.

\begin{thebibliography}{9}

\bibitem{Lagarias1985}
J. C. Lagarias,
``The $3x+1$ problem and its generalizations,''
\emph{The American Mathematical Monthly}, vol. 92, no. 1, pp. 3--23, 1985.
doi: \href{https://doi.org/10.1080/00029890.1985.11971528}{10.1080/00029890.1985.11971528}.

\bibitem{MatthewsWatts1985}
K. Matthews and A. Watts,
``A Markov approach to the generalized Syracuse algorithm,''
\emph{Acta Arithmetica}, vol. 45, no. 1, pp. 29--42, 1985.
doi: \href{https://doi.org/10.4064/aa-45-1-29-42}{10.4064/aa-45-1-29-42}.

\bibitem{Carnielli2015}
W. Carnielli,
``Some natural generalizations of the Collatz problem,''
\emph{Applied Mathematics E-Notes}, vol. 15, pp. 207--215, 2015.

\bibitem{BanazadehBase2026}
F. Banazadeh,
``Ternary $(5k\pm1)/4$ Collatz-Type Map with Two Observed Attractors:
Exhaustive Computational Verification up to $10^{10}$ and Algebraic Structure,''
Zenodo, 2026. doi: \href{https://doi.org/10.5281/zenodo.22709065}{10.5281/zenodo.22709065}.

\bibitem{BanazadehScaled2026}
F. Banazadeh,
``A $4^r$-Scaled Family of the Ternary $(4k-1(2))/3$ Collatz-Type Map:
Exact Algebraic Conjugacy, Cycle Preservation, and Transfer of the $10^9$ Finite Verification,''
Zenodo, 2026. doi: \href{https://doi.org/10.5281/zenodo.22671949}{10.5281/zenodo.22671949}.

\end{thebibliography}

\end{document}
